\section{The no-arbitrage band}\label{sec:band}

\Cref{prop:band} says that the basis should behave differently inside and outside a band: inside, it follows the
demand of outside traders and nothing pulls it back; outside, arbitrage pushes it toward the band. This section
estimates where the band lies, using days on which the stock is not in a ban.

\subsection{A threshold model}
Let $x_{t-1}$ be yesterday's basis minus the stock's median basis on non-ban days, in basis points, and $\Delta m_t$ the
change in the basis of the same contract from yesterday to today. Measuring the basis relative to each stock's own median
removes permanent differences between stocks, such as a stock that always trades slightly rich. The three-regime
threshold model is
\begin{equation}\label{eq:setar}
\Delta m_t=a_r+\rho_r\,x_{t-1}+e_t,\qquad
r=\begin{cases}\text{lower} & x_{t-1}<g_L,\\ \text{middle} & g_L\le x_{t-1}\le g_U,\\ \text{upper} & x_{t-1}>g_U.\end{cases}
\end{equation}
The thresholds $g_L$ and $g_U$ are the estimated edges of the band, and $\rho_r$ measures how strongly the basis moves
back in each regime: $\rho_r=0$ means no pull at all and $\rho_r=-1$ means that a deviation disappears within a day.
The model predicts $\rho$ close to zero in the middle regime and clearly negative outside it. I choose the two
thresholds by least squares over a grid of percentiles of $x$, keeping at least 10\% of the observations in each outer
regime and 20\% in the middle. Because the thresholds are chosen from the data, the usual F-test against a linear model
is not valid; I compare the largest F-statistic over the grid with its distribution under the linear model, obtained
by a wild bootstrap clustered by trading day. Confidence intervals for the thresholds come from resampling whole
stocks.

\subsection{Estimates}
The sample has 473,110 stock-days for 360 stocks. Linearity is rejected decisively: the test statistic is 7,471 against
a 95\% critical value of 31. The band runs from $-34.5$\bp\ to $+23.9$\bp\ around each stock's median, about 58\bp\
wide, and both edges are estimated precisely (\cref{tab:setar}). The basis lies inside the band on 76\% of days and in
each outer regime on 12\%. As the model predicts, the band is asymmetric: the lower edge lies further from the centre
than the upper edge. \Cref{fig:bandscatter} shows the pattern directly: the further the basis is from the stock's
median yesterday, the more it moves back today, with the strongest pull above the upper edge.

\begin{table}[t]
\caption{The basis moves in a wide, asymmetric band, but much of the daily reversion is noise}\label{tab:setar}
\exsub{Pooled three-regime threshold model, non-ban days, 2015 to 2026 (473,110 observations, 360 stocks)}
\small\begin{tabular}{@{}>{\raggedright\arraybackslash\bfseries\color{accent}}p{6.4cm}ccc@{}}
\toprule
& \thd{Lower regime} & \thd{Middle regime} & \thd{Upper regime} \\
\midrule
Threshold (bp, relative to the stock's median) & below $-34.5$ & & above $+23.9$ \\
95\% interval (stock bootstrap) & $[-37.2,\,-32.2]$ & & $[+20.3,\,+26.1]$ \\
Share of observations & 0.12 & 0.76 & 0.12 \\
Slope $\rho_r$ & $-0.084$ & $-0.468$ & $-0.809$ \\
Slope, yesterday's deviation instrumented & $-0.121$ & $-0.192$ & $-0.141$ \\
\bottomrule
\end{tabular}
\exsource{Thresholds chosen by least squares over a grid of percentiles; intervals from resampling whole stocks. In the last row
yesterday's deviation is instrumented by the deviation two days earlier, and regimes are defined on that deviation. Source:
author's calculations.}
\end{table}

\begin{figure}[t]
\caption{The further the basis moves from the stock's median, the more it moves back the next day}\label{fig:bandscatter}
\exsub{Mean change in the basis today by yesterday's deviation from the stock's median (5\bp\ bins), non-ban days; shaded: the estimated band}
\centering\includegraphics[width=\linewidth]{fig_band_scatter.pdf}
\exsource{Source: author's calculations.}
\end{figure}

\subsection{How much of the reversion is noise}
The slopes need care. Inside the band the model predicts almost no pull, yet the middle slope of $-0.47$ implies that
half of a deviation disappears within about a day. The likely reason is measurement noise. The daily volume-weighted
prices of the future and the stock are averages over slightly different trades, so the measured basis contains an error
that does not persist. An error in yesterday's basis appears in $x_{t-1}$ with one sign and in $\Delta m_t$ with the
other, which pushes the estimated slope toward $-1$ even if the true basis does not move back at all. A standard remedy
is to instrument yesterday's deviation with the deviation two days earlier, whose error is unrelated to yesterday's. With
this correction all three slopes become small and similar ($-0.12$, $-0.19$ and $-0.14$, last row of \cref{tab:setar}).
Much of the day-to-day reversion is therefore noise, and the faster pull above the upper edge in the uncorrected model is
not a robust finding. The thresholds, which describe where the basis spends its time rather than how fast it returns, are
the more reliable part of the estimates.

\subsection{Smooth adjustment}
Arbitrageurs differ in their costs, so the market's response may strengthen gradually rather than switch on at one
threshold (Tse, 2001). I therefore also estimate smooth-transition models of the form
\begin{equation}\label{eq:star}
\Delta m_t=a+\big[\rho_1+\rho_2\,G(x_{t-1})\big]x_{t-1}+e_t ,
\end{equation}
where $G$ moves smoothly from 0 to 1. With an exponential $G$, which is symmetric around a centre, deviations in either
direction are treated alike; with a logistic $G$ the pull can differ on the two sides. By the Bayesian information
criterion the threshold model fits best, the logistic model almost as well (with weak reversion below about $+5$\bp\
and strong reversion above), and the symmetric exponential model clearly worse. A symmetric model does not suit an
asymmetric band.

\subsection{The band over time and over the life of a contract}
Estimated year by year, the band is stable (\cref{fig:bandyears}). The exception is 2020, when the market dislocation
of March pushed the lower edge to $-59$\bp. The band also depends on the time left to expiry, as \cref{prop:cs} (iii)
predicts: the lower edge is at $-40$\bp\ when 17 or more trading days remain and at $-28$\bp\ in the last ten days. In a
smooth-transition model in which the strength of reversion depends on the days left to expiry, reversion is stronger in
the last week of a contract.

\begin{figure}[t]
\caption{The band is stable from year to year, apart from the dislocation of March 2020}\label{fig:bandyears}
\exsub{Thresholds of the pooled threshold model by calendar year, 2015 to 2025, non-ban days, basis relative to each stock's median}
\centering\includegraphics[width=\linewidth]{fig_band_years.pdf}
\exsource{Years with a change in the tax are examined in \cref{sec:costs}. Source: author's calculations.}
\end{figure}

\subsection{Bands of individual stocks}
Estimating the model separately for each of the 258 stocks with at least 500 non-ban days gives a median band of $-26$ to
$+17$\bp, about 41\bp\ wide, with large differences between stocks: the middle 80\% of stocks have widths between
about 18 and 76\bp. These bands are much wider than the intraday bands of earlier Indian studies (\cref{tab:litbands}), which is
expected: an arbitrageur who holds a position overnight pays the taxes on delivery and bears the risk of carrying it,
while an intraday trader does neither. \Cref{sec:costs} asks which costs and frictions explain the location of the band
and its differences across stocks.
